\unitchapter{1}{8}{Information \& Communication Technology (ICT)}{p1-08}

\unitmeta{Expected questions: \textasciitilde5 (10 marks) · Target: 5/5 ---
easiest unit for a CS candidate}

\begin{sylbox}

\begin{itemize}
\tightlist
\item[$\square$]
  ICT: general abbreviations and terminology
\item[$\square$]
  Basics of Internet, Intranet, E-mail, Audio and Video-conferencing
\item[$\square$]
  Digital initiatives in higher education
\item[$\square$]
  ICT and Governance
\end{itemize}

\end{sylbox}

\needspace{125pt}

\hypertarget{p1-08--1-computer-basics}{%
\section{Computer Basics}\label{p1-08--1-computer-basics}}

For Paper I you need a clear picture of what the parts of a computer do rather than how they are built.
Every computer, from a phone to a supercomputer, follows the same plan: data enter through input devices,
are processed by the CPU using instructions held in memory, and leave through output devices.

\begin{figure}[!htbp]
\centering
\begin{adjustbox}{max width=\linewidth}
\begin{tikzpicture}
  \node[slate, text width=24mm] (IN) at (0,0) {\textbf{Input}\\\footnotesize keyboard, mouse, scanner};
  \node[box, fill=fgrey, minimum width=46mm, minimum height=34mm] (CPU) at (5.2,0) {};
  \node[capt, anchor=north] at (CPU.north) {central processing unit};
  \node[sand, minimum width=36mm] (CU) at (5.2,0.45) {Control unit};
  \node[sand, minimum width=17mm, font=\footnotesize] (ALU) at (4.27,-0.75) {ALU};
  \node[sand, minimum width=17mm, font=\footnotesize] (REG) at (6.13,-0.75) {Registers};
  \node[sage, text width=24mm] (OUT) at (10.4,0) {\textbf{Output}\\\footnotesize monitor, printer};
  \node[rose, text width=36mm] (MEM) at (5.2,-3.1) {\textbf{Primary memory}\\\footnotesize RAM, ROM, cache};
  \node[db, minimum width=30mm, minimum height=13mm, text width=26mm] (SEC) at (10.4,-3.1) {\textbf{Secondary}\\\footnotesize HDD, SSD, pen drive};
  \draw[arr] (IN) -- (CPU); \draw[arr] (CPU) -- (OUT);
  \draw[arrd] (CPU) -- (MEM); \draw[arrd] (MEM) -- (SEC);
\end{tikzpicture}
\end{adjustbox}
\caption{The basic organisation of a computer. The CPU fetches instructions and data from primary memory;
secondary storage holds them permanently.}
\end{figure}


\needspace{125pt}

\hypertarget{p1-08--memory-hierarchy-fastcostly-at-top}{%
\subsection{Memory hierarchy (fast/costly at top)}\label{p1-08--memory-hierarchy-fastcostly-at-top}}

No single memory technology is both fast and cheap, so computers combine several in a hierarchy. Frequently
used data are kept in the small, fast levels near the processor; everything else waits in the large, slow
levels below.

\begin{figure}[!htbp]
\centering
\begin{tikzpicture}
  \foreach \i/\w/\name/\col in {0/9.0/{Optical disc / tape}/fgrey, 1/7.4/{SSD / hard disk}/fslate, 2/5.8/{Main memory (DRAM)}/fsage, 3/4.2/{Cache (SRAM)}/fsand, 4/2.6/{Registers}/frose} {
    \pic at ({(9-\w)/2},{\i*0.9}) {slab={\w}{0.9}{1.0}{\col}};
    \node[font=\small] at (4.5,{\i*0.9+0.45}) {\name};
  }
  \draw[arr] (10.6,0) -- node[lbl, rotate=90, below] {faster, costlier, smaller} (10.6,4.5);
\end{tikzpicture}
\caption{The memory hierarchy. Moving up, each level is faster and more expensive per byte but smaller.}
\end{figure}


\needspace{95pt}

\hypertarget{p1-08--units}{%
\subsection{Units}\label{p1-08--units}}

\begin{center}
\small
\renewcommand{\arraystretch}{1.3}
\begin{tabular}{l l @{\hspace{2.2em}} l l @{\hspace{2.2em}} l l}
\toprule
1 nibble & 4 bits & 1 KB & $2^{10}$ B & 1 PB & $2^{50}$ B \\
1 byte & 8 bits & 1 MB & $2^{20}$ B & 1 EB & $2^{60}$ B \\
 & & 1 GB & $2^{30}$ B & 1 ZB & $2^{70}$ B \\
 & & 1 TB & $2^{40}$ B & 1 YB & $2^{80}$ B \\
\bottomrule
\end{tabular}
\end{center}


\needspace{125pt}

\hypertarget{p1-08--number-systems-frequent-question}{%
\subsection{Number Systems (frequent question!)}\label{p1-08--number-systems-frequent-question}}

Computers store everything in binary because electronic circuits have two reliable states. Octal and
hexadecimal are compact ways of writing binary: each octal digit stands for three bits and each hexadecimal
digit for four.

\begin{itemize}
\tightlist
\item
  Binary (2), Octal (8), Decimal (10), Hexadecimal (16).
\item
  (25)₁₀ = (11001)₂ = (31)₈ = (19)₁₆
\item
  (1010.101)₂ = 10.625
\item
  Group 3 bits → octal, 4 bits → hex.
\end{itemize}

\needspace{130pt}

\hypertarget{p1-08--2-abbreviations-must-know}{%
\section{Abbreviations (must know)}\label{p1-08--2-abbreviations-must-know}}

\begingroup\small

\begin{longtable}[]{@{}
  >{\raggedright\arraybackslash}p{(\columnwidth - 2\tabcolsep) * \real{0.2089}}
  >{\raggedright\arraybackslash}p{(\columnwidth - 2\tabcolsep) * \real{0.7911}}@{}}
\toprule\noalign{}
\begin{minipage}[b]{\linewidth}\raggedright
\textbf{Abbr.}
\end{minipage} & \begin{minipage}[b]{\linewidth}\raggedright
\textbf{Full form}
\end{minipage} \\
\midrule\noalign{}
\endhead
\bottomrule\noalign{}
\endlastfoot
ALU & Arithmetic Logic Unit \\
BIOS & Basic Input Output System \\
CMOS & Complementary Metal Oxide Semiconductor \\
DNS & Domain Name System \\
DHCP & Dynamic Host Configuration Protocol \\
FTP & File Transfer Protocol \\
HTML & HyperText Markup Language \\
HTTP(S) & HyperText Transfer Protocol (Secure) \\
IMAP & Internet Message Access Protocol \\
ISP & Internet Service Provider \\
LAN/MAN/WAN & Local/\allowbreak{}Metropolitan/\allowbreak{}Wide Area Network \\
MODEM & Modulator--Demodulator \\
OCR / OMR / MICR & Optical Character / Optical Mark / Magnetic Ink Character Recognition \\
POP3 & Post Office Protocol v3 \\
SMTP & Simple Mail Transfer Protocol \\
TCP/IP & Transmission Control Protocol / Internet Protocol \\
URL & Uniform Resource Locator \\
VoIP & Voice over Internet Protocol \\
VPN & Virtual Private Network \\
Wi-Fi & Wireless Fidelity \\
GIF, JPEG, PNG & Graphics Interchange Format, Joint Photographic Experts Group, Portable Network Graphics \\
PDF & Portable Document Format \\
CC / BCC & Carbon Copy / Blind Carbon Copy \\
IoT & Internet of Things \\
LMS & Learning Management System (Moodle) \\
NIC & National Informatics Centre \\
CERT-In & Indian Computer Emergency Response Team \\
\end{longtable}

\endgroup

\needspace{125pt}

\hypertarget{p1-08--3-internet-intranet-extranet}{%
\section{Internet, Intranet, Extranet}\label{p1-08--3-internet-intranet-extranet}}

The same Internet technologies --- TCP/IP, web browsers, e-mail --- can be used on a private network
(intranet), shared with trusted partners (extranet) or opened to the world (Internet). What differs is who
is allowed in.

\begin{figure}[!htbp]
\centering
\begin{tikzpicture}
  \node[panel, minimum width=58mm, minimum height=36mm] (ORG) at (0,0) {};
  \node[capt, anchor=north west] at (ORG.north west) {the organisation};
  \node[slate, text width=30mm] (IN) at (0,-0.1) {\textbf{Intranet}\\\footnotesize private; employees only};
  \node[sage, text width=30mm] (EX) at (0,-3.4) {\textbf{Extranet}\\\footnotesize intranet opened to selected partners};
  \node[brick, minimum width=6mm, minimum height=26mm, inner sep=0pt] (FW) at (4.0,0) {};
  \foreach \y in {-1.0,-0.5,...,1.0} \draw[fgink!60, line width=0.3pt] (3.7,\y) -- (4.3,\y);
  \node[lbl, below=1mm of FW] {firewall};
  \node[cloudnode, minimum width=34mm, minimum height=20mm, fill=fgrey] (NET) at (7.4,0) {\textbf{Internet}\\\footnotesize public, global};
  \draw[arrd] (IN) -- (FW); \draw[arrd] (FW) -- (NET);
  \draw[dasharr, <->] (IN) -- (EX);
\end{tikzpicture}
\caption{Intranet, extranet and Internet. A firewall filters traffic between the private network and the
public Internet.}
\end{figure}


\begin{itemize}
\tightlist
\item
  \textbf{Internet} originated from \textbf{ARPANET (1969)}, US DoD. \textbf{WWW} invented by \textbf{Tim Berners-Lee (1989, CERN)}.
\item
  India\textquotesingle s first: \textbf{ERNET (1986)}; public internet by \textbf{VSNL on 15 Aug 1995}.
\item
  IPv4 = 32 bits; IPv6 = 128 bits. Domain types: .com, .org, .edu, .gov, .ac.in, .nic.in.
\item
  Browser = client software (Chrome, Firefox); Search engine = website (Google, Bing).
\item
  \textbf{Cookies} = small data files stored by websites in the browser.
\end{itemize}

\needspace{125pt}

\hypertarget{p1-08--4-e-mail}{%
\section{E-mail}\label{p1-08--4-e-mail}}

E-mail is a store-and-forward system: messages wait on servers until the recipient collects them, which is
why sender and recipient need not be online at the same time.

\begin{figure}[!htbp]
\centering
\begin{tikzpicture}[node distance=24mm]
  \node[slate, text width=22mm] (A) {sender's\\mail client};
  \node[db, text width=20mm, right=of A] (S1) {sender's\\server};
  \node[db, text width=20mm, right=of S1] (S2) {receiver's\\server};
  \node[sage, text width=22mm, right=of S2] (B) {receiver's\\mail client};
  \draw[arr] (A) -- node[lbl, above] {SMTP} node[lbl, below] {send} (S1);
  \draw[arr] (S1) -- node[lbl, above] {SMTP} node[lbl, below] {relay} (S2);
  \draw[arr] (B) -- node[lbl, above] {POP3 / IMAP} node[lbl, below] {retrieve} (S2);
\end{tikzpicture}
\caption{How e-mail travels. SMTP \emph{pushes} mail towards the recipient's server; POP3 or IMAP lets
the recipient \emph{pull} it down.}
\end{figure}


\begin{itemize}
\tightlist
\item
  \textbf{SMTP} = sending (port 25/587); \textbf{POP3} = download \& delete (110); \textbf{IMAP} = sync, mail stays on server (143).
\item
  First email: \textbf{Ray Tomlinson (1971)}, introduced \texttt{@}.
\item
  \textbf{BCC} recipients are hidden from others. \textbf{Spam} = unsolicited mail. \textbf{Phishing} = fraudulent mail to steal data.
\item
  Attachments: MIME (Multipurpose Internet Mail Extensions).
\end{itemize}

\needspace{95pt}

\hypertarget{p1-08--5-audio--video-conferencing}{%
\section{Audio \& Video Conferencing}\label{p1-08--5-audio--video-conferencing}}

\begin{itemize}
\tightlist
\item
  Real-time communication over IP: Zoom, Google Meet, Microsoft Teams, Webex; government: \textbf{JioMeet, Bharat VC (NIC\textquotesingle s)}.
\item
  Protocols: \textbf{H.323}, \textbf{SIP} (Session Initiation Protocol), RTP. Codecs compress audio/video.
\item
  Synchronous (live) vs asynchronous (recorded, forums, email) learning.
\item
  Webinar = web seminar; Podcast = audio episodes; Vodcast = video podcast.
\end{itemize}

\needspace{160pt}

\hypertarget{p1-08--6-digital-initiatives-in-higher-education}{%
\section{Digital Initiatives in Higher Education}\label{p1-08--6-digital-initiatives-in-higher-education}}

These initiatives, coordinated largely by the Ministry of Education, UGC, AICTE and INFLIBNET, together
form India\textquotesingle s digital infrastructure for higher education. Learn each name with its purpose.

\begingroup\small

\begin{longtable}[]{@{}
  >{\raggedright\arraybackslash}p{(\columnwidth - 2\tabcolsep) * \real{0.4711}}
  >{\raggedright\arraybackslash}p{(\columnwidth - 2\tabcolsep) * \real{0.5289}}@{}}
\toprule\noalign{}
\begin{minipage}[b]{\linewidth}\raggedright
\textbf{Initiative}
\end{minipage} & \begin{minipage}[b]{\linewidth}\raggedright
\textbf{Purpose}
\end{minipage} \\
\midrule\noalign{}
\endhead
\bottomrule\noalign{}
\endlastfoot
SWAYAM & MOOCs (2017) \\
SWAYAM PRABHA & DTH educational TV channels \\
NDL India & National Digital Library (IIT Kharagpur) \\
e-PG Pathshala & PG e-content (UGC/\allowbreak{}INFLIBNET) \\
Shodhganga / Shodhgangotri & Theses / synopses repository \\
e-ShodhSindhu & E-journal consortium (merged UGC-INFONET, N-LIST, INDEST) \\
\textbf{ONOS (One Nation One Subscription)} & Nation-wide access to scholarly journals (2024/25) \\
NAD / DigiLocker & National Academic Depository --- digital certificates \\
\textbf{ABC} & Academic Bank of Credits (NEP 2020) \\
\textbf{NIRF} & National Institutional Ranking Framework (2015) \\
AISHE & All India Survey on Higher Education \\
Virtual Labs, Spoken Tutorial, e-Yantra, FOSSEE & Lab/skill initiatives \\
SAMARTH & e-Governance suite for universities \\
VIDWAN & Database of experts (INFLIBNET) \\
PM e-VIDYA & Multi-mode access to digital education (2020) \\
NEAT & National Educational Alliance for Technology (AICTE, AI-based adaptive learning) \\
DIKSHA & Teachers\textquotesingle{} platform \\
\end{longtable}

\endgroup

\needspace{125pt}

\hypertarget{p1-08--7-ict--governance}{%
\section{ICT \& Governance}\label{p1-08--7-ict--governance}}

E-governance uses ICT to deliver government services more quickly, transparently and cheaply. It is
usually described by the parties it connects.

\begin{figure}[!htbp]
\centering
\begin{adjustbox}{max width=\linewidth}
\begin{tikzpicture}
  \node[circ, minimum size=24mm, fill=fwine!15, font=\small\bfseries, align=center] (G) at (0,0) {e-Gover-\\nance};
  \node[slate, text width=44mm] (a) at (-4.3,1.4) {\textbf{G2C} — Government to Citizen\\\footnotesize UMANG, DigiLocker, Passport Seva};
  \node[sage, text width=44mm] (b) at (4.3,1.4) {\textbf{G2B} — Government to Business\\\footnotesize GeM, GSTN, MCA21};
  \node[sand, text width=44mm] (c) at (-4.3,-1.4) {\textbf{G2G} — Government to Government\\\footnotesize e-Office, PFMS};
  \node[rose, text width=44mm] (d) at (4.3,-1.4) {\textbf{G2E} — Government to Employee\\\footnotesize e-HRMS};
  \foreach \n in {a,b,c,d} \draw[thinline] (G) -- (\n);
\end{tikzpicture}
\end{adjustbox}
\caption{The four relationships served by e-governance, with representative Indian platforms.}
\end{figure}


\begin{itemize}
\tightlist
\item
  \textbf{Digital India} (1 July 2015) --- 9 pillars incl. broadband highways, e-Kranti, Information for All.
\item
  Key platforms: \textbf{UMANG} (unified app), \textbf{DigiLocker}, \textbf{GeM} (Government e-Marketplace), \textbf{BHIM-UPI} (NPCI), \textbf{CoWIN}, \textbf{MyGov}, \textbf{e-Hospital}, \textbf{Bhashini} (language AI).
\item
  \textbf{IT Act 2000} (amended 2008) --- legal recognition of e-records \& digital signatures; Sec 66A struck down (Shreya Singhal, 2015). \textbf{DPDP Act 2023}.
\item
  Cyber security: firewall, antivirus, encryption, 2FA. Malware: virus (needs host), worm (self-replicating), Trojan (disguised), ransomware, spyware.
\end{itemize}

\needspace{166pt}

\hypertarget{p1-08--8-deeper-dive--generations-software-file-types--emerging-tech}{%
\section{Deeper Dive --- Generations, Software, File Types \& Emerging Tech}\label{p1-08--8-deeper-dive--generations-software-file-types--emerging-tech}}

\needspace{130pt}

\hypertarget{p1-08--81-generations-of-computers}{%
\subsection{Generations of Computers}\label{p1-08--81-generations-of-computers}}

\begingroup\small

\begin{longtable}[]{@{}
  >{\raggedright\arraybackslash}p{(\columnwidth - 6\tabcolsep) * \real{0.1259}}
  >{\raggedright\arraybackslash}p{(\columnwidth - 6\tabcolsep) * \real{0.1673}}
  >{\raggedright\arraybackslash}p{(\columnwidth - 6\tabcolsep) * \real{0.3783}}
  >{\raggedright\arraybackslash}p{(\columnwidth - 6\tabcolsep) * \real{0.3286}}@{}}
\toprule\noalign{}
\begin{minipage}[b]{\linewidth}\raggedright
\textbf{Generation}
\end{minipage} & \begin{minipage}[b]{\linewidth}\raggedright
\textbf{Period}
\end{minipage} & \begin{minipage}[b]{\linewidth}\raggedright
\textbf{Technology}
\end{minipage} & \begin{minipage}[b]{\linewidth}\raggedright
\textbf{Example}
\end{minipage} \\
\midrule\noalign{}
\endhead
\bottomrule\noalign{}
\endlastfoot
1st & 1940--56 & Vacuum tubes; machine language & ENIAC, UNIVAC \\
2nd & 1956--63 & Transistors; assembly, FORTRAN, COBOL & IBM 1401 \\
3rd & 1964--71 & Integrated circuits; OS, multiprogramming & IBM 360 \\
4th & 1971--present & Microprocessors (VLSI); PCs, GUI & Intel 4004 (first microprocessor) \\
5th & Present \& beyond & AI, ULSI, parallel processing, quantum & --- \\
\end{longtable}

\endgroup

India: \textbf{PARAM 8000 (1991)} by C-DAC (first Indian supercomputer); \textbf{National Supercomputing Mission (2015)}; AIRAWAT (AI supercomputer).

\needspace{125pt}

\hypertarget{p1-08--82-software}{%
\subsection{Software}\label{p1-08--82-software}}

Hardware does nothing without instructions. Software is commonly divided into the programs that run the
computer itself and the programs that do useful work for the user.

\begin{figure}[!htbp]
\centering
\begin{adjustbox}{max width=\linewidth}
\begin{tikzpicture}
  \node[grey, minimum width=30mm] (S) at (0,0) {\textbf{Software}};
  \node[slate, minimum width=40mm] (SY) at (-3.6,-1.4) {\textbf{System software}};
  \node[sage, minimum width=40mm] (AP) at (3.6,-1.4) {\textbf{Application software}};
  \draw[arr] (S) -- (SY); \draw[arr] (S) -- (AP);
  \node[plain, text width=50mm, align=left, font=\footnotesize, anchor=north] at (-3.6,-2.05)
    {operating systems — Windows, Linux, macOS, Android\\utilities — antivirus, disk clean-up, compression\\translators — compiler, interpreter, assembler\\device drivers and firmware};
  \node[plain, text width=50mm, align=left, font=\footnotesize, anchor=north] at (3.6,-2.05)
    {general purpose — word processor, spreadsheet, browser\\special purpose — accounting (Tally), payroll, CAD\\\phantom{x}\\\phantom{x}};
\end{tikzpicture}
\end{adjustbox}
\caption{System software manages the machine; application software serves the user's tasks.}
\end{figure}


\begin{itemize}
\tightlist
\item
  \textbf{Open source} (source available, free to modify: Linux, LibreOffice, Moodle) vs \textbf{proprietary} (MS Office).
\item
  \textbf{Freeware} (free, closed), \textbf{shareware} (trial), \textbf{firmware} (software in ROM/flash).
\end{itemize}

\needspace{130pt}

\hypertarget{p1-08--83-file-extensions--office-shortcuts}{%
\subsection{File Extensions \& Office Shortcuts}\label{p1-08--83-file-extensions--office-shortcuts}}

\begingroup\small

\begin{longtable}[]{@{}
  >{\raggedright\arraybackslash}p{(\columnwidth - 6\tabcolsep) * \real{0.1458}}
  >{\raggedright\arraybackslash}p{(\columnwidth - 6\tabcolsep) * \real{0.2202}}
  >{\raggedright\arraybackslash}p{(\columnwidth - 6\tabcolsep) * \real{0.1202}}
  >{\raggedright\arraybackslash}p{(\columnwidth - 6\tabcolsep) * \real{0.2399}}@{}}
\toprule\noalign{}
\begin{minipage}[b]{\linewidth}\raggedright
\textbf{Extension}
\end{minipage} & \begin{minipage}[b]{\linewidth}\raggedright
\textbf{Type}
\end{minipage} & \begin{minipage}[b]{\linewidth}\raggedright
\textbf{Extension}
\end{minipage} & \begin{minipage}[b]{\linewidth}\raggedright
\textbf{Type}
\end{minipage} \\
\midrule\noalign{}
\endhead
\bottomrule\noalign{}
\endlastfoot
.docx & Word document & .xlsx & Spreadsheet \\
.pptx & Presentation & .pdf & Portable document \\
.txt & Plain text & .csv & Comma-separated values \\
.jpg/\allowbreak{}.png/\allowbreak{}.gif & Images & .mp3/.wav & Audio \\
.mp4/.avi & Video & .zip/.rar & Compressed \\
.exe & Executable (Windows) & .html & Web page \\
\end{longtable}

\endgroup

\begingroup\small

\begin{longtable}[]{@{}
  >{\raggedright\arraybackslash}p{(\columnwidth - 6\tabcolsep) * \real{0.1732}}
  >{\raggedright\arraybackslash}p{(\columnwidth - 6\tabcolsep) * \real{0.1976}}
  >{\raggedright\arraybackslash}p{(\columnwidth - 6\tabcolsep) * \real{0.1345}}
  >{\raggedright\arraybackslash}p{(\columnwidth - 6\tabcolsep) * \real{0.1911}}@{}}
\toprule\noalign{}
\begin{minipage}[b]{\linewidth}\raggedright
\textbf{Shortcut}
\end{minipage} & \begin{minipage}[b]{\linewidth}\raggedright
\textbf{Action}
\end{minipage} & \begin{minipage}[b]{\linewidth}\raggedright
\textbf{Shortcut}
\end{minipage} & \begin{minipage}[b]{\linewidth}\raggedright
\textbf{Action}
\end{minipage} \\
\midrule\noalign{}
\endhead
\bottomrule\noalign{}
\endlastfoot
Ctrl + C / X / V & Copy / cut / paste & Ctrl + Z / Y & Undo / redo \\
Ctrl + S & Save & Ctrl + P & Print \\
Ctrl + F & Find & Ctrl + H & Replace \\
Ctrl + A & Select all & F7 & Spell check (Word) \\
F5 & Slideshow / refresh & Alt + F4 & Close window \\
\end{longtable}

\endgroup

Spreadsheet: cell address = column letter + row number (B3); \textbf{absolute reference} \$B\$3; functions SUM, AVERAGE, COUNT, IF, VLOOKUP.

\needspace{130pt}

\hypertarget{p1-08--84-emerging-technologies-frequently-asked-one-liners}{%
\subsection{Emerging Technologies (frequently asked one-liners)}\label{p1-08--84-emerging-technologies-frequently-asked-one-liners}}

\begingroup\small

\begin{longtable}[]{@{}
  >{\raggedright\arraybackslash}p{(\columnwidth - 2\tabcolsep) * \real{0.3089}}
  >{\raggedright\arraybackslash}p{(\columnwidth - 2\tabcolsep) * \real{0.6617}}@{}}
\toprule\noalign{}
\begin{minipage}[b]{\linewidth}\raggedright
\textbf{Term}
\end{minipage} & \begin{minipage}[b]{\linewidth}\raggedright
\textbf{Meaning}
\end{minipage} \\
\midrule\noalign{}
\endhead
\bottomrule\noalign{}
\endlastfoot
Cloud computing & On-demand computing over the internet: SaaS, PaaS, IaaS \\
IoT & Physical devices with sensors connected to the internet \\
Artificial Intelligence & Machines performing tasks needing human intelligence; ChatGPT, Bhashini \\
Machine learning & Systems that learn patterns from data \\
Blockchain & Distributed, tamper-evident ledger (Bitcoin) \\
Big data & Volume, velocity, variety (+ veracity, value) \\
AR / VR & Augmented reality overlays digital on real; virtual reality is fully immersive \\
5G & Fifth-generation mobile; launched in India on 1 Oct 2022 \\
Quantum computing & Qubits, superposition; India\textquotesingle s National Quantum Mission (2023) \\
Open educational resources (OER) & Free, openly licensed learning materials (Creative Commons) \\
\end{longtable}

\endgroup

\needspace{95pt}

\hypertarget{p1-08--85-internet-basics-extended}{%
\subsection{Internet Basics Extended}\label{p1-08--85-internet-basics-extended}}

\begin{itemize}
\tightlist
\item
  Web 1.0 (read-only) → Web 2.0 (read-write, social media) → Web 3.0 (semantic, decentralised).
\item
  \textbf{Search operators}: "exact phrase", \texttt{site:ac.in}, \texttt{filetype:pdf}, minus sign to exclude.
\item
  \textbf{Bandwidth} measured in bps; \textbf{latency} in ms. Wired (Ethernet, fibre) vs wireless (Wi-Fi, 4G/5G, satellite).
\item
  \textbf{IP address} identifies a device; \textbf{MAC address} identifies a network card (48 bits).
\end{itemize}

\needspace{95pt}

\hypertarget{p1-08--previous-year-questions-pyq-pattern}{%
\section{Previous Year Questions (PYQ pattern)}\label{p1-08--previous-year-questions-pyq-pattern}}

\begin{enumerate}
\def\labelenumi{\arabic{enumi}.}
\tightlist
\item
  (1011)₂ + (0111)₂ = ? → \textbf{(10010)₂}
\item
  (237)₈ in decimal = 2×64+3×8+7 = \textbf{159}
\item
  1 TB equals: \textbf{1024 GB}
\item
  Which protocol is used to send e-mail? \textbf{SMTP}
\item
  Which protocol keeps mail on the server and syncs folders? \textbf{IMAP}
\item
  BCC stands for: \textbf{Blind Carbon Copy}
\item
  Which is NOT an input device? (a) scanner (b) light pen (c) plotter (d) OMR --- \textbf{Ans: (c)}
\item
  WWW was invented by: \textbf{Tim Berners-Lee}
\item
  The internet originated from: \textbf{ARPANET}
\item
  Private network inside an organisation using internet technology: \textbf{Intranet}
\item
  URL stands for: \textbf{Uniform Resource Locator}
\item
  Which is a volatile memory? \textbf{RAM}
\item
  Which converts domain names to IP addresses? \textbf{DNS}
\item
  IPv6 address length: \textbf{128 bits}
\item
  "Academic Bank of Credits" is an initiative under: \textbf{NEP 2020}
\item
  The portal providing Indian theses: \textbf{Shodhganga}
\item
  The national ranking framework for HEIs: \textbf{NIRF}
\item
  Which is a G2B service? \textbf{GeM / GSTN}
\item
  Malware that replicates itself without a host program: \textbf{Worm}
\item
  Fraudulent attempt to obtain credentials via email: \textbf{Phishing}
\item
  Video conferencing protocol: \textbf{H.323 / SIP}
\item
  The Information Technology Act was passed in: \textbf{2000}
\item
  Arrange in increasing order of capacity: \textbf{KB \textless{} MB \textless{} GB \textless{} TB \textless{} PB}
\end{enumerate}

\textbf{More practice questions}

\begin{enumerate}
\def\labelenumi{\arabic{enumi}.}
\setcounter{enumi}{23}
\tightlist
\item
  The first Indian supercomputer: \textbf{PARAM 8000}
\item
  Integrated circuits were used in: \textbf{third-generation computers}
\item
  Which is open-source software? (a) MS Word (b) LibreOffice Writer (c) Photoshop (d) Windows --- \textbf{Ans: (b)}
\item
  Absolute cell reference in a spreadsheet: \textbf{\$A\$1}
\item
  Shortcut to undo: \textbf{Ctrl + Z}
\item
  A .csv file stores: \textbf{tabular data as comma-separated text}
\item
  Web 2.0 is characterised by: \textbf{user-generated content and interactivity}
\item
  Software stored permanently in ROM: \textbf{firmware}
\item
  Which is NOT system software? (a) OS (b) compiler (c) spreadsheet (d) device driver --- \textbf{Ans: (c)}
\item
  5G services were launched in India in: \textbf{October 2022}
\item
  Search operator to restrict results to a website: \textbf{site:}
\end{enumerate}

\begin{revbox}

\begin{itemize}
\tightlist
\item
  SMTP send · POP3 download · IMAP sync
\item
  ARPANET 1969 · WWW 1989 (Berners-Lee) · India public internet 1995 (VSNL)
\item
  IT Act 2000 · Digital India 2015 · DPDP Act 2023
\end{itemize}

\end{revbox}
